\documentclass[11pt,a4paper]{article}
\usepackage[utf8]{inputenc}
\usepackage{geometry}
\geometry{margin=1in}
\usepackage{hyperref}
\usepackage{booktabs}
\usepackage{enumitem}
\usepackage{listings}
\usepackage{amsmath}
\usepackage{xcolor}

\lstset{
  language=C++,
  basicstyle=\ttfamily\small,
  keywordstyle=\bfseries,
  commentstyle=\itshape,
  showstringspaces=false,
  breaklines=true,
  captionpos=b,
  numbers=left,
  numberstyle=\tiny\color{gray},
}

\title{Digital Design Training Program \\ \large Day 3: Multiplexing and Hardware Timers}
\author{Authored by }
\date{April 1, 2026}

\begin{document}

\maketitle

\section*{Theme: Overcoming pin limitations by wiring a display matrix and using visual persistence, while managing time with hardware interrupts.}

\subsection*{Module 3.1: The Multiplexing Matrix (Hardware) (60 Minutes)}

\begin{itemize}
    \item \textbf{Goal:} Understand how to control many components (6 displays) using a minimal number of microcontroller pins by sharing data lines and individually switching power.
    \item \textbf{Requirements:} 6x Common-Anode 7-Segment Displays, 1x 7447 IC, 6x PNP Transistors (e.g., 2N3906), 6x 1k$\Omega$ Resistors, Breadboards, Jumper Wires.
    \item \textbf{Conceptual Overview (The Railroad Analogy):} Imagine you have one train (the BCD data from the 7447) and 6 different stations (the displays). Instead of building 6 separate tracks, you build one main track that connects to all stations. To send the train to a specific station, you first set the switches (the transistors) to direct the train to, say, Station 3, and then you send the train. Then you quickly reset the switches to direct to Station 5 and send the next train. If you do this fast enough, it looks like every station is getting its own personal train service simultaneously.
    \item \textbf{Wiring Guide (The Display Bus):}
\end{itemize}

\begin{table}[h]
\centering
\begin{tabular}{lp{4cm}p{6.5cm}}
\toprule
\textbf{Component} & \textbf{Connection} & \textbf{Notes} \\
\midrule
Displays 1-6 & All segment `a' pins tied together, all 'b' pins together, etc. & This creates a shared 7-wire ``data bus''. All displays receive the same segment data at the same time. \\
Shared Bus & 7447 Outputs (a-g) via 220$\Omega$ resistors & The 7447 determines \emph{what} number pattern is sent down the data bus. \\
PNP Transistors & Emitter to 5V, Collector to one display's Common Anode & The transistor acts as a digital switch for the display's power supply. \\
Arduino A0-A5 & Each pin connects to a transistor's Base via a 1k$\Omega$ resistor & The Arduino now has control over which power switch is ON, determining \emph{which} display is allowed to light up. \\
\bottomrule
\end{tabular}
\end{table}

\begin{itemize}
    \item \textbf{Action Steps:}
    \begin{enumerate}[label=\arabic*.]
        \item \textbf{The Transistor as a Switch:} We use a PNP transistor. When its Base pin is at the same voltage as its Emitter (5V), it's OFF. When we use the Arduino to pull its Base pin LOW (to 0V), current flows from Emitter to Collector, turning the display's power ON. The 1k$\Omega$ resistor on the base is crucial to limit the current drawn from the Arduino pin.
        \item Assemble the hardware matrix carefully. This is the most complex wiring so far. Use different colored wires for the data bus and the control lines to keep things organized.
    \end{enumerate}
    \item \textbf{Checkpoint:} The hardware matrix is complete. You should be able to trace the path of the 7 data lines to all displays, and the 6 separate control lines from the Arduino to each transistor.
\end{itemize}
\newpage
\subsection*{Module 3.2: Software Multiplexing (The Fast Loop) (45 Minutes)}

\begin{itemize}
    \item \textbf{Goal:} Exploit a biological phenomenon called Persistence of Vision (POV) by rapidly flashing digits one by one, making them appear as a single, stable display.
    \item \textbf{Conceptual Overview:} Your eyes and brain retain an image for a fraction of a second after it disappears. If we flash digit 1, then digit 2, then digit 3, etc., faster than about 24 times per second, your brain merges them into one continuous image. We will turn one display on, show a number, turn it off, then move to the next.
    \item \textbf{Action Steps \& Code:}
\begin{lstlisting}[caption={Code for Module 3.2: Software Multiplexing}]
// Pins for the 7447 BCD data bus
const int bcd_a = 8;
const int bcd_b = 9;
const int bcd_c = 10;
const int bcd_d = 11;

// Array of pins to control the PNP transistor switches
const int digit_pins[] = {A0, A1, A2, A3, A4, A5};
const int N_DIGITS = 6;

// The data we want to display. Let's start with ``123456''.
int time_arr[] = {1, 2, 3, 4, 5, 6};

// Function to send a number (0-9) to the 7447
void show_bcd(int num) {
  digitalWrite(bcd_a, (num >> 0) & 1);
  digitalWrite(bcd_b, (num >> 1) & 1);
  digitalWrite(bcd_c, (num >> 2) & 1);
  digitalWrite(bcd_d, (num >> 3) & 1);
}

void setup() {
  // Set BCD pins as outputs
  pinMode(bcd_a, OUTPUT);
  pinMode(bcd_b, OUTPUT);
  pinMode(bcd_c, OUTPUT);
  pinMode(bcd_d, OUTPUT);

  // Set digit control pins as outputs and initialize them to HIGH.
  // HIGH turns the PNP transistors OFF, so no displays are lit initially.
  for (int i = 0; i < N_DIGITS; i++) {
    pinMode(digit_pins[i], OUTPUT);
    digitalWrite(digit_pins[i], HIGH);
  }
}

void loop() {
  // This is the core multiplexing loop.
  for (int i = 0; i < N_DIGITS; i++) {
    // 1. Set the data bus for the number we want to show.
    show_bcd(time_arr[i]);
    
    // 2. Enable ONLY the desired display by pulling its transistor base LOW.
    digitalWrite(digit_pins[i], LOW);
    
    // 3. Wait a very short time. This is how long the digit is visible.
    delay(2); // 2 milliseconds
    
    // 4. IMPORTANT: Turn the display OFF before changing the data.
    // This prevents ``ghosting'' where segments from the next digit
    // briefly appear on the current display.
    digitalWrite(digit_pins[i], HIGH);
  }
}
\end{lstlisting}
    \item \textbf{Checkpoint:} The display shows ``123456'' clearly and without flickering.
    \item \textbf{Challenge:} Change the \verb|delay(2);| to \verb|delay(100);|. You should now be able to see the numbers scanning across the display, one at a time. This proves the multiplexing concept. Now, set it back to 2. What happens if you remove the final \verb|digitalWrite(digit_pins[i], HIGH);|? You should see ghosting.
\end{itemize}
\newpage
\subsection*{Module 3.3: Bare-Metal Timers (Timer1 ISR) (60 Minutes)}

\begin{itemize}
    \item \textbf{Goal:} Replace the blocking \verb|delay()| function with a non-blocking hardware timer interrupt to create a precise, background timekeeping ``tick''.
    \item \textbf{Conceptual Overview:} The \verb|delay()| function is simple but terrible for responsive programs. It completely halts the microcontroller. For our clock, the multiplexing in \verb|loop()| must run as fast as possible, while our timekeeping needs to happen at a steady, predictable rate. We achieve this with a \textbf{hardware timer interrupt}. We will configure one of the Arduino's built-in timers to send a signal (an ``interrupt'') to the CPU at a precise frequency. The CPU will pause its main work (the \verb|loop()|), run a special function called an Interrupt Service Routine (ISR), and then resume the loop.
    \item \textbf{Action Steps \& Code:}
\begin{lstlisting}[caption={Code for Module 3.3: 10Hz Hardware Interrupt}]
// This variable is shared between the ISR and other code.
// 'volatile' tells the compiler that this value can change at any time,
// preventing it from making optimizations that could break the code.
volatile int deciseconds = 0;
const int led_pin = 13; // Onboard LED for visual feedback

void setup() {
  pinMode(led_pin, OUTPUT);

  // --- Timer1 Configuration ---
  // We are directly manipulating the microcontroller's registers.
  
  cli(); // Disable all interrupts globally while we configure.

  // TCCR1A & TCCR1B are the Timer/Counter Control Registers for Timer1.
  TCCR1A = 0; // Set TCCR1A to all zeros. We don't need PWM modes.
  TCCR1B = 0; // Set TCCR1B to all zeros initially.

  // Set the Compare Match Register (OCR1A). This is the value the
  // timer will count up to.
  // FORMULA: OCR1A = (Clock Speed / (Prescaler * Target Freq)) - 1
  // (16,000,000 / (256 * 10Hz)) - 1 = 6249
  OCR1A = 6249;

  // Configure TCCR1B:
  // WGM12 = 1: Sets CTC (Clear Timer on Compare Match) mode.
  // CS12 = 1: Sets the prescaler to 256.
  TCCR1B |= (1 << WGM12) | (1 << CS12);

  // TIMSK1 is the Timer Interrupt Mask Register.
  // OCIE1A = 1: Enables the interrupt for when the timer matches OCR1A.
  TIMSK1 |= (1 << OCIE1A);

  sei(); // Re-enable all interrupts globally.
}

// This is the Interrupt Service Routine.
// The name TIMER1_COMPA_vect is special; it's defined by the hardware.
// This function is AUTOMATICALLY called by the CPU every time the
// Timer1 counter matches our OCR1A value (i.e., 10 times per second).
ISR(TIMER1_COMPA_vect) {
  deciseconds++; // Increment our timekeeping variable.
  
  // Toggle the LED so we can see the interrupt is working.
  // At 10Hz, the LED will change state 10 times/sec, or blink 5 times/sec.
  digitalWrite(led_pin, !digitalRead(led_pin));
}

void loop() {
  // The loop is now empty!
  // All the multiplexing code would go here, and it would run
  // completely uninterrupted by timekeeping duties.
  // The ISR handles time in the background.
}
\end{lstlisting}
    \item \textbf{Checkpoint:} The onboard LED is blinking steadily at 5Hz. This confirms your hardware timer is running correctly and firing the ISR at exactly 10Hz, completely independent of the (currently empty) \verb|loop()| function.
    \item \textbf{Challenge:} Change the interrupt frequency to 1Hz. You will need to change the prescaler and/or the OCR1A value. A prescaler of 1024 is also available (CS12 and CS10 bits). Calculate the new OCR1A value for a 1Hz tick using a 1024 prescaler.
\end{itemize}

\end{document}
